\section{Kesimpulan}
Milestone 2 berhasil merealisasikan komponen parser untuk Arion Compiler dengan 
algoritma Recursive Descent yang terstruktur. Parser menerima daftar token dari 
lexer atau token stream, memverifikasi kesesuaian token dengan grammar Arion, 
serta membangun parse tree sebagai representasi hierarkis program. Implementasi 
yang modular memisahkan tanggung jawab antara pembacaan input, pemrosesan token, 
parsing, dan pencetakan parse tree, sehingga alur kerja mudah ditelusuri.

Hasil pengujian dirancang untuk menunjukkan bahwa parser mampu menangani konstruksi 
sintaks utama seperti deklarasi, pernyataan kontrol, dan ekspresi, serta melaporkan 
kesalahan sintaks secara informatif dengan lokasi baris dan kolom. Dengan demikian, 
parser Milestone 2 telah menyediakan fondasi yang diperlukan untuk melanjutkan ke 
tahap semantic analysis pada milestone berikutnya.

\section{Saran}
Pengembangan pada milestone selanjutnya dapat memanfaatkan parse tree yang sudah 
dihasilkan sebagai masukan untuk semantic analysis, termasuk pengecekan tipe, 
scope, dan konsistensi deklarasi. Selain itu, pengujian dapat diperluas dengan 
menambahkan variasi kasus yang lebih kompleks, terutama pada struktur subprogram, 
record, dan array multi-dimensi agar cakupan grammar tervalidasi lebih luas.

Ke depan, kualitas error handling juga dapat ditingkatkan melalui strategi 
recovery sederhana agar parser mampu melanjutkan parsing setelah menemukan 
error pertama. Hal ini akan membuat pelaporan error lebih lengkap dalam satu 
kali eksekusi dan mempermudah proses debugging bagi pengguna.
